% !PRINTED none
% Part-title / section-opening page: no folio printed. Centred divider "Zwölfter
% Abschnitt. / Ordnungen im quadratischen Körper" with a footnote marker, a rule,
% then the "§ 96. Diskriminanten der Ordnungen." heading and the start of running
% text (Definition 1, remarks 1.-2., footnote). Sequence implies printed page 351.

\begin{center}
\gesperrt{Zwölfter Abschnitt.}

Ordnungen im quadratischen Körper\textsuperscript{1}).
\end{center}

\rule{3cm}{0.4pt}

\begin{center}
\textbf{\S\ 96. \gesperrt{Diskriminanten der Ordnungen.}}
\end{center}

1. \gesperrt{Definition. Ist $Q$ eine natürliche Zahl, so heißt die Gesamtheit der ganzen Zahlen des quadratischen Körpers $\Omega$, die nach dem Modul $Q$ mit einer rationalen Zahl kongruent sind, eine Ordnung, und $Q$ heißt der Führer dieser Ordnung.} Die Ordnung mit dem Führer $Q$ wird mit $[Q]$ bezeichnet. Es ergibt sich aus dieser Definition:

1. \gesperrt{Alle rationalen ganzen Zahlen gehören jeder Ordnung an;}

2. \gesperrt{Summe, Differenz und Produkt zweier Zahlen der Ordnung gehören derselben Ordnung an;}

d.\,h. man kann innerhalb der Ordnung Addition, Subtraktion und Multiplikation unbedingt ausführen, nicht aber allgemein die Division.

Die Gesamtheit aller \gesperrt{ganzen Zahlen} des Körpers $\Omega$ bilde gleichfalls eine Ordnung (vom Führer 1), die \gesperrt{Hauptordnung}, die also mit $[1]$ zu bezeichnen wäre.

\rule{4cm}{0.4pt}

\textsuperscript{1}) \gesperrt{Dedekind} hat den Ausdruck Ordnung in der allgemeinen Theorie der algebraischen Zahlen eingeführt, weil sie in dem besonderen Fall des quadratischen Körpers den \gesperrt{Gaussschen} Ordnungen der quadratischen Formen entsprechen. \gesperrt{Hilbert} gebraucht dafür den Ausdruck \glqq Ring\grqq{} oder \glqq Zahlring\grqq{}. Vgl. \gesperrt{Dirichlet-Dedekind}, Vorlesungen über Zahlentheorie (\S\ 172 der dritten, \S\ 170 der vierten Auflage). \gesperrt{Dedekind}: Über die Anzahl der Idealklassen in den verschiedenen Ordnungen eines endlichen Körpers, Festschrift zur Säkularfeier von \gesperrt{Gauss'} Geburtstag, Braunschweig 1877. \gesperrt{Hilbert}: Die Theorie der algebraischen Zahlen, Bericht der Deutschen Mathematischen Vereinigung von 1894/95, Kap. IX.
